\documentclass[11pt,a4paper]{scrartcl}

\usepackage{amssymb}
\usepackage{amsmath}

\usepackage[utf8]{inputenc}
\usepackage[T1]{fontenc}
\newcommand{\ats}[1]{\par\noindent\textit{Atsakymas:} #1}
\usepackage{cancel}
\usepackage{tikz}
\usepackage{lmodern} 
\usepackage{amsmath,amssymb,amsthm}
\usepackage{graphicx}
\usepackage[dvipsnames,svgnames]{xcolor}
\usepackage{hyperref}
\usepackage{mathrsfs}
\usepackage[shortlabels]{enumitem}
\usepackage[obeyFinal,textsize=scriptsize,shadow,loadshadowlibrary]{todonotes}
\usepackage{mathtools}
\usepackage{microtype}

%%% Paragraph layout
\setlength{\parindent}{0pt}
\setlength{\parskip}{0.6\baselineskip}

%%% Colored sections
\renewcommand*{\sectionformat}%
  {\color{purple}\S\thesection\enskip}
\renewcommand*{\subsectionformat}%
  {\color{purple}\S\thesubsection\enskip}
\renewcommand*{\subsubsectionformat}%
  {\color{purple}\S\thesubsubsection\enskip}
\addtokomafont{paragraph}{\color{orange!35!black}\P\ }
\KOMAoptions{numbers=noenddot}

%%% Title
\addtokomafont{subtitle}{\Large}
\setkomafont{author}{\Large\scshape}
\setkomafont{date}{\Large\normalsize}
% Neigiamas vertikalus tarpas, kuris pakelia dokumento antraštę į viršų. 
% Galite keisti -2.5cm į -3cm (aukščiau) arba -1.5cm (žemiau).
\titlehead{\vspace*{-7.5cm}} 

% PRIDĖTA: Automatiškai perrašome datą į mokyklos pavadinimą
\AtBeginDocument{%
  \date{\normalsize Vilniaus licėjus}
}

%%% Page setup
\usepackage[headsepline]{scrlayer-scrpage}
\renewcommand{\headfont}{}
\addtolength{\textheight}{3.14cm}
\setlength{\footskip}{0.5in}
\setlength{\headsep}{10pt}
% Puslapio antraštė turi rodyti tik konkrečius dokumento duomenis.
% Nustatome juos aiškiai, kad neatsirastų "\@author" / "\@title" arba senų reikšmių.
\makeatletter
\AtBeginDocument{%
  \ihead{\footnotesize\textbf{Andrius Berniukevičius}}%
  \ohead{\footnotesize\textbf{\@title}}%
}
\makeatother
\automark{section}
\chead{}
\cfoot{\pagemark}

%%% Macros
\providecommand{\ol}{\overline}
\providecommand{\ul}{\underline}
\providecommand{\wt}{\widetilde}
\providecommand{\wh}{\widehat}
\providecommand{\eps}{\varepsilon}
\providecommand{\half}{\frac{1}{2}}
\providecommand{\inv}{^{-1}}
\newcommand{\dang}{\measuredangle} %% Directed angle
\newcommand{\vocab}[1]{\textbf{\color{blue}\sffamily #1}}
\providecommand{\CC}{\mathbb C}
\providecommand{\FF}{\mathbb F}
\providecommand{\NN}{\mathbb N}
\providecommand{\QQ}{\mathbb Q}
\providecommand{\RR}{\mathbb R}
\providecommand{\ZZ}{\mathbb Z}
\providecommand{\ts}{\textsuperscript}
\providecommand{\dg}{^\circ}
\providecommand{\ii}{\item}
\providecommand{\defeq}{\coloneqq}
\DeclareMathOperator*{\lcm}{lcm}
\DeclareMathOperator*{\argmin}{arg min}
\DeclareMathOperator*{\argmax}{arg max}
\providecommand{\hrulebar}{\par
\hspace{\fill}\rule{0.95\linewidth}{.7pt}\hspace{\fill}
\par\nointerlineskip \vspace{\baselineskip}}

%%% Optional colored theorems
\usepackage{thmtools}
\usepackage[framemethod=TikZ]{mdframed}
\mdfdefinestyle{mdbluebox}{%
  roundcorner=10pt,
  linewidth=1pt,
  skipabove=12pt,
  innerbottommargin=9pt,
  skipbelow=2pt,
  linecolor=blue,
  nobreak=true,
  backgroundcolor=TealBlue!5,
}
\declaretheoremstyle[
  headfont=\sffamily\bfseries\color{MidnightBlue},
  mdframed={style=mdbluebox},
  headpunct={\\[3pt]},
  postheadspace={0pt}
]{thmbluebox}
\mdfdefinestyle{mdredbox}{%
  linewidth=0.5pt,
  skipabove=12pt,
  frametitleaboveskip=5pt,
  frametitlebelowskip=0pt,
  skipbelow=2pt,
  frametitlefont=\bfseries,
  innertopmargin=4pt,
  innerbottommargin=8pt,
  nobreak=true,
  backgroundcolor=Salmon!5,
  linecolor=RawSienna,
}
\declaretheoremstyle[
  headfont=\bfseries\color{RawSienna},
  mdframed={style=mdredbox},
  headpunct={\\[3pt]},
  postheadspace={0pt},
]{thmredbox}
\mdfdefinestyle{mdgreenbox}{%
  skipabove=8pt,
  linewidth=2pt,
  rightline=false,
  leftline=true,
  topline=false,
  bottomline=false,
  linecolor=ForestGreen,
  backgroundcolor=ForestGreen!5,
}
\declaretheoremstyle[
  headfont=\bfseries\sffamily\color{ForestGreen!70!black},
  bodyfont=\normalfont,
  spaceabove=2pt,
  spacebelow=1pt,
  mdframed={style=mdgreenbox},
  headpunct={ --- },
]{thmgreenbox}
\mdfdefinestyle{mdblackbox}{%
  skipabove=8pt,
  linewidth=3pt,
  rightline=false,
  leftline=true,
  topline=false,
  bottomline=false,
  linecolor=black,
  backgroundcolor=RedViolet!5!gray!5,
}
\declaretheoremstyle[
  headfont=\bfseries,
  bodyfont=\normalfont\small,
  spaceabove=0pt,
  spacebelow=0pt,
  mdframed={style=mdblackbox}
]{thmblackbox}

%% Aplinkos su rėmeliais (thmtools)
\declaretheorem[style=thmbluebox,name=Teorema]{theorem}
\declaretheorem[style=thmbluebox,name=Lema]{lemma}
\declaretheorem[style=thmbluebox,name=Savybė]{property}
\declaretheorem[style=thmbluebox,name=Teiginys]{proposition}
\declaretheorem[style=thmbluebox,name=Išvada]{corollary}
\declaretheorem[style=thmbluebox,name=Teorema,numbered=no]{theorem*}
\declaretheorem[style=thmbluebox,name=Lema,numbered=no]{lemma*}
\declaretheorem[style=thmbluebox,name=Teiginys,numbered=no]{proposition*}
\declaretheorem[style=thmbluebox,name=Išvada,numbered=no]{corollary*}
\declaretheorem[style=thmgreenbox,name=Algoritmas]{algorithm}
\declaretheorem[style=thmgreenbox,name=Algoritmas,numbered=no]{algorithm*}
\declaretheorem[style=thmgreenbox,name=Teiginys]{claim}
\declaretheorem[style=thmgreenbox,name=Teiginys,numbered=no]{claim*}
\declaretheorem[style=thmredbox,name=Pavyzdys]{example}
\declaretheorem[style=thmredbox,name=Pavyzdys,numbered=no]{example*}
\declaretheorem[style=thmblackbox,name=Pastaba]{remark}
\declaretheorem[style=thmblackbox,name=Pastaba,numbered=no]{remark*}

%% Standartinės amsthm aplinkos (Apibrėžimai ir užduotys)
\theoremstyle{definition}
\ifcsname conjecture\endcsname\else\newtheorem{conjecture}{Hipotezė}\fi
\ifcsname definition\endcsname\else\newtheorem{definition}{Apibrėžimas}\fi
\ifcsname fact\endcsname\else\newtheorem{fact}{Faktas}\fi
\ifcsname answer\endcsname\else\newtheorem{answer}{Atsakymas}\fi
\ifcsname ques\endcsname\else\newtheorem{ques}{Klausimas}\fi
\ifcsname exercise\endcsname\else\newtheorem{exercise}{Užduotis}\fi
\ifcsname problem\endcsname\else\newtheorem{problem}{Uždavinys}[section]\fi
\ifcsname conjecture*\endcsname\else\newtheorem*{conjecture*}{Hipotezė}\fi
\ifcsname definition*\endcsname\else\newtheorem*{definition*}{Apibrėžimas}\fi
\ifcsname fact*\endcsname\else\newtheorem*{fact*}{Faktas}\fi
\ifcsname answer*\endcsname\else\newtheorem*{answer*}{Atsakymas}\fi
\ifcsname ques*\endcsname\else\newtheorem*{ques*}{Klausimas}\fi
\ifcsname exercise*\endcsname\else\newtheorem*{exercise*}{Užduotis}\fi
\ifcsname problem*\endcsname\else\newtheorem*{problem*}{Uždavinys}\fi

\newcounter{answerssection}
\newcommand{\answerssection}{%
  \setcounter{answerssection}{\value{section}}%
  \section{Atsakymai}%
  \setcounter{problem}{0}%
  \renewcommand{\theproblem}{\arabic{answerssection}.\arabic{problem}}%
}

%% GALUTINIS SPRENDIMAS PROOF APLINKAI
%% --- PRADŽIA: Įrodymo aplinkos sutvarkymas ---
\makeatletter

% 1. Užtikriname, kad žodis būtų "Įrodymas", net jei "babel" paketas bando jį perrašyti
\AtBeginDocument{%
  \renewcommand{\proofname}{Įrodymas}%
  \@ifpackageloaded{babel}{%
    \addto\captionslithuanian{\renewcommand{\proofname}{Įrodymas}}%
  }{}%
}

% 2. Priverstinai pašaliname scrartcl klasės bold šriftą ir paliekame TIK italic
% 3. Sujungiame su tuščiu kvadratėliu dešiniajame kampe.
\renewcommand{\qedsymbol}{\ensuremath{\Box}}
\renewenvironment{proof}[1][\proofname]{\par
  \pushQED{\qed}%
  \normalfont \topsep6\p@\@plus6\p@\relax
  \trivlist
  \item[\hskip\labelsep
        \normalfont\itshape % Priverčia naudoti tik pasvirą šriftą, kaip nuotraukoje
    #1\@addpunct{.}]\ignorespaces
}{%
  \popQED\endtrivlist\@endpefalse
}
\newenvironment{solution}{\par
  \normalfont \topsep6\p@\@plus6\p@\relax
  \trivlist
  \item[\hskip\labelsep
        \normalfont\itshape
    \textit{Sprendimas}\@addpunct{.}]\ignorespaces
}{%
  \endtrivlist\@endpefalse
}
\newcommand{\ans}[1]{\par\noindent\textit{Atsakymas:} #1}
\makeatother


\title{Skaičių teorija: DBD, MBK ir Euklido algoritmas}
\author{Andrius Berniukevičius}

\newenvironment{solution}
    {\begin{list}{}{\setlength{\leftmargin}{10pt}\setlength{\rightmargin}{10pt}}\item[]}
    {\end{list}}

\begin{document}

\maketitle

\section{Bendrųjų daliklių ir kartotinių paieška}

Matematikoje, sprendžiant lygtis su sveikaisiais skaičiais, suprastinant trupmenas ar ieškant pasikartojančių ciklų, nuolat susiduriame su dviem esminiais dydžiais:
\begin{itemize}
    \item \textbf{DBD (Didžiausias bendras daliklis)} – didžiausias natūralusis skaičius, iš kurio tobulai dalijasi abu duotieji skaičiai.
    \item \textbf{MBK (Mažiausias bendras kartotinis)} – mažiausias natūralusis skaičius, kuris tobulai dalijasi iš abiejų duotųjų skaičių.
\end{itemize}

Praeitame užsiėmime išmokome kiekvieną natūralųjį skaičių išskaidyti pirminiais dauginamaisiais (užrašyti jo \textbf{kanoninį skaidinį}). Žinant skaičių kanoninius skaidinius, rasti jų DBD ir MBK yra labai paprasta.

\begin{center}
\fbox{
\begin{minipage}{0.9\textwidth}
\textbf{Kanoninio skaidinio metodas:}

Tarkime, turime du skaičius ir jų skaidymus pirminiais dauginamaisiais:
$A = 2^3 \cdot 3^2 \cdot 5^1$ (tai yra $360$) ir $B = 2^2 \cdot 3^4$ (tai yra $324$).

\textbf{Kaip rasti DBD?} Iš abiejų skaidinių išsirenkame tik tuos pirminius dauginamuosius, kuriuos turi \textbf{abu} skaičiai, ir imame \textbf{mažiausią} jų laipsnį (nes daliklis turi tilpti į abu skaičius). \\
$DBD(360, 324) = 2^2 \cdot 3^2 = 4 \cdot 9 = \mathbf{36}$.

\textbf{Kaip rasti MBK?} Surenkame \textbf{visus} paminėtus pirminius dauginamuosius ir imame \textbf{didžiausią} jų laipsnį (kad gautas skaičius dalintųsi iš abiejų pradinių skaičių). \\
$MBK(360, 324) = 2^3 \cdot 3^4 \cdot 5^1 = 8 \cdot 81 \cdot 5 = \mathbf{3240}$.
\end{minipage}
}
\end{center}

% ==========================================
% KAI METODAS PAVEDA IR EUKLIDO ALGORITMAS
% ==========================================
\section{Kai skaidymas pirminiais buksuoja: Euklido algoritmas}
Kanoninio skaidinio metodas yra patogus, kai skaičiai nedideli. Bet įsivaizduokite, kad olimpiadoje gaunate užduotį: \textit{Raskite $DBD(1073, 851)$}. 
Bandyti juos išskaidyti pirminiais dauginamaisiais užtruktų labai ilgai. Kaip rasti atsakymą greitai?

Čia į pagalbą ateina daugiau nei prieš 2000 metų senovės Graikijoje aprašytas \textbf{Euklido algoritmas}. Jis remiasi fundamentalia dalumo savybe:
\[ \mathbf{DBD(a, b) = DBD(a - b, b)} \]
\textit{(Atėmus mažesnį skaičių iš didesnio, jų didžiausias bendras daliklis nepasikeičia).}

\begin{center}
\fbox{
\begin{minipage}{0.9\textwidth}
\textbf{Kodėl Euklido algoritmas veikia? (Griežtas įrodymas)}

Įrodykime, kad skaičių poros $(a, b)$ ir poros $(a-b, b)$ bendri dalikliai yra visiškai tie patys.
\begin{itemize}
    \item \textbf{Pirma kryptis:} Tarkime, skaičius $d$ yra bendras skaičių $a$ ir $b$ daliklis. Tai reiškia, kad $d \mid a$ ir $d \mid b$. Pagal mums jau žinomą Skirtumo taisyklę (iš pirmojo užsiėmimo), $d$ privalo dalinti ir jų skirtumą: $d \mid (a-b)$. Vadinasi, $d$ yra ir skaičių $(a-b)$ bei $b$ bendras daliklis.
    \item \textbf{Antra kryptis:} Tarkime, kažkoks skaičius $k$ dalija $(a-b)$ ir dalija $b$. Tai reiškia, $k \mid (a-b)$ ir $k \mid b$. Pagal Sumos taisyklę, $k$ privalo dalinti jų sumą: $k \mid ((a-b) + b)$, kas reiškia, jog $k \mid a$. Vadinasi, $k$ yra ir pradinių skaičių $a$ bei $b$ bendras daliklis.
\end{itemize}
\textbf{Išvada:} Kadangi abiejų porų visų įmanomų daliklių sąrašai yra visiškai identiški, tai ir patys \textbf{didžiausi} tų sąrašų skaičiai privalo sutapti! Todėl $DBD(a, b) = DBD(a-b, b)$.
\end{minipage}
}
\end{center}

\textbf{Euklido algoritmo taikymas (dalijimas su liekana):}
Raskime $DBD(1073, 851)$.
Vietoj to, kad skaidytume, pasinaudokime įrodyta taisykle – atimkime!
$1073 - 851 = 222$. Vadinasi:
$DBD(1073, 851) = DBD(851, 222)$.

Skaičiai sumažėjo. Tęsiame procesą – atimame 222 iš 851 tiek kartų, kiek telpa (tai yra tas pats, kas \textbf{dalijimas su liekana}).
Skaičius $222$ į $851$ telpa $3$ kartus ($222 \cdot 3 = 666$). Atimame: $851 - 666 = 185$.
Gauname: $DBD(851, 222) = DBD(222, 185)$.

Tęsiame: $222 - 185 = 37$.
Gauname: $DBD(222, 185) = DBD(185, 37)$.

Kiek kartų $37$ telpa į $185$? Lygiai $5$ kartus ($37 \cdot 5 = 185$). Liekana lygi nuliui!
Kai gauname liekaną nulis, procesas baigtas. Paskutinis nelygus nuliui skaičius ir yra ieškomas DBD.
Atsakymas: \textbf{DBD yra 37}.
% ==========================================
% AUKSINĖ TAISYKLĖ
% ==========================================
% ==========================================
% AUKSINĖ TAISYKLĖ
% ==========================================

\section{Tarpusavyje pirminiai skaičiai ir MBK formulė}

O kaip rasti MBK dideliems skaičiams? Egzistuoja griežta algebrinė formulė, siejanti abu dydžius:
\[ \mathbf{DBD(a, b) \cdot MBK(a, b) = a \cdot b} \]
\textit{Dviejų skaičių DBD ir MBK sandauga visada yra lygi pačių tų skaičių sandaugai.}
\textbf{Įrodymas:}\\
Prisiminkime, kaip ieškome DBD ir MBK iš skaičių kanoninių skaidinių. Tarkime, turime bet kokį pirminį skaičių $p$, kuris į skaičių $a$ įeina laipsniu $x$, o į skaičių $b$ – laipsniu $y$. 
(Pavyzdžiui, jei $a = 2^3 \cdot \dots$ ir $b = 2^5 \cdot \dots$, tai $x=3$, o $y=5$).

\begin{itemize}
    \item Ieškodami \textbf{DBD}, mes pasirenkame \textbf{mažiausią} iš šių dviejų laipsnių. Pažymėkime jį $\min(x, y)$.
    \item Ieškodami \textbf{MBK}, mes pasirenkame \textbf{didžiausią} laipsnį. Pažymėkime jį $\max(x, y)$.
\end{itemize}

Dauginant skaičius, vienodų pagrindų laipsniai yra sudedami.
Pažiūrėkime, koks bus pirminio skaičiaus $p$ laipsnis abiejose mūsų lygties pusėse:
1) Sandaugoje $DBD(a,b) \cdot MBK(a,b)$ laipsnis bus: $\min(x, y) + \max(x, y)$.
2) Sandaugoje $a \cdot b$ laipsnis bus tiesiog: $x + y$.

Bet juk akivaizdu, kad paėmus du skaičius (pvz., $3$ ir $5$), mažesniojo ir didesniojo suma yra lygiai ta pati kaip ir pačių skaičių suma ($3 + 5 = 8$). Matematiškai: 
\[ \min(x, y) + \max(x, y) = x + y \]
Kadangi ši lygybė galioja absoliučiai kiekvienam pirminiam dauginamajam $p$, abiejų lygties pusių kanoniniai skaidiniai yra identiški. Vadinasi, $DBD \cdot MBK = a \cdot b$. \hfill $\square$

\vspace{0.5cm}

\textbf{Tarpusavyje pirminiai skaičiai:} 
Jeigu dviejų skaičių $DBD(a, b) = 1$ (jie neturi jokių bendrų pirminių dauginamųjų, išskyrus 1), tokie skaičiai vadinami \textbf{tarpusavyje pirminiais}.
Iš ką tik įrodytos MBK formulės išplaukia, kad tarpusavyje pirminių skaičių mažiausias bendras kartotinis yra lygus tiesiog jų sandaugai ($1 \cdot MBK = a \cdot b$).


% PAVYZDŽIAI
% ==========================================
\section{Olimpiadiniai pavyzdžiai}

\begin{example}[Algebrinis Euklido algoritmo taikymas]
Įrodykite, kad bet kokiam natūraliajam skaičiui $n$, trupmena $\frac{7n+3}{5n+2}$ yra nesuprastinama.
\end{example}

\begin{solution}
\textbf{Sprendimas:}\\
Trupmena yra nesuprastinama tada ir tik tada, kai jos skaitiklis ir vardiklis yra tarpusavyje pirminiai skaičiai (jų DBD lygus 1). Pritaikysime Euklido algoritmą:
$DBD(7n+3, 5n+2)$.

Atimame mažesnį reiškinį iš didesnio: $(7n+3) - (5n+2) = 2n+1$.
Vadinasi, $DBD(7n+3, 5n+2) = DBD(5n+2, 2n+1)$.

Tęsiame atiminėjimą. Kiek kartų $2n+1$ telpa į $5n+2$? Telpa 2 kartus ($4n+2$). Atimame:
$(5n+2) - 2(2n+1) = (5n+2) - (4n+2) = n$.
Vadinasi, $DBD(5n+2, 2n+1) = DBD(2n+1, n)$.

Atimame toliau: $(2n+1) - 2 \cdot n = 1$.
Vadinasi, $DBD(2n+1, n) = DBD(n, 1)$.
Didžiausias bendras skaičiaus $n$ ir vieneto daliklis yra \textbf{1}. 
Kadangi $DBD = 1$, skaitiklis ir vardiklis yra tarpusavyje pirminiai, todėl trupmena nesuprastinama jokiam $n$. \hfill $\square$
\end{solution}

\begin{example}[Skaičių atstatymas]
Raskite natūraliuosius skaičius $a$ ir $b$, jeigu žinoma, kad $DBD(a, b) = 15$, o jų sandauga $a \cdot b = 6750$.
\end{example}

\begin{solution}
\textbf{Sprendimas:}\\
Kadangi $DBD(a, b) = 15$, abu skaičiai dalijasi iš $15$. 
Užrašykime juos algebriškai: $a = 15x$ ir $b = 15y$. 
Čia kintamieji $x$ ir $y$ privalo būti \textbf{tarpusavyje pirminiai} ($DBD(x, y) = 1$), kitaip pradinis DBD būtų buvęs didesnis nei 15.

Įstatykime į duotą sandaugą:
\[ 15x \cdot 15y = 6750 \]
\[ 225 \cdot xy = 6750 \]
\[ xy = 30 \]

Ieškome dviejų tarpusavyje pirminių skaičių, kurių sandauga lygi 30.
Galimi variantai porai $(x, y)$:
1) $1$ ir $30$.  Tuomet $a = 15 \cdot 1 = 15$, o $b = 15 \cdot 30 = 450$.
2) $2$ ir $15$.  Tuomet $a = 15 \cdot 2 = 30$, o $b = 15 \cdot 15 = 225$.
3) $3$ ir $10$.  Tuomet $a = 15 \cdot 3 = 45$, o $b = 15 \cdot 10 = 150$.
4) $5$ ir $6$.   Tuomet $a = 15 \cdot 5 = 75$, o $b = 15 \cdot 6 = 90$.

Visi šie variantai tenkina sąlygas. \hfill $\square$
\end{solution}

\vspace{0.5cm}

% ==========================================
% UŽDAVINIAI
% ==========================================
\newpage
\section{Uždaviniai savarankiškam sprendimui (16 iššūkių)}

\begin{enumerate}
    \renewcommand{\labelenumi}{\textbf{\theenumi.}}
    \item Naudodamiesi kanoninio skaidinio metodu, raskite $DBD$ ir $MBK$ skaičiams $360$ ir $504$.
    \item Naudodamiesi Euklido algoritmu (dalijimu su liekana), raskite $DBD(3132, 7200)$.
    \item Raskite $DBD(2026, 2027)$. Ką galite pasakyti apie bet kokių dviejų iš eilės einančių natūraliųjų skaičių $n$ ir $n+1$ bendrus daliklius?
    \item Žinoma, kad ieškome $DBD(n, n+2)$. Kokias reikšmes gali įgyti šis dydis? Pagrįskite atsakymą Euklido algoritmu.
    \item Dviejų skaičių $a$ ir $b$ didžiausias bendras daliklis yra $8$, o jų mažiausias bendras kartotinis yra $144$. Kokia yra šių dviejų skaičių sandauga $a \cdot b$?
    \item Euklido algoritmu įrodykite, kad trupmena $\frac{3n+1}{2n+1}$ yra nesuprastinama.
    \item Euklido algoritmu įrodykite, kad trupmena $\frac{12n+5}{5n+2}$ yra nesuprastinama jokiam natūraliajam $n$.
    \item Žinoma, kad dviejų skaičių sandauga yra $1500$, o jų DBD yra $10$. Raskite visus tokius skaičius $a$ ir $b$.
    \item Dviejų natūraliųjų skaičių suma lygi $168$, o jų $DBD$ yra $24$. Raskite visus tokius skaičius.
    \item Raskite mažiausią natūralųjį skaičių, kuris dalijant iš $2$ duoda liekaną $1$, dalijant iš $3$ duoda liekaną $2$, iš $4$ duoda $3$, iš $5$ duoda $4$, o iš $6$ duoda $5$.
    \item Koks gali būti didžiausias bendras daliklis $DBD(a+b, a-b)$, jeigu žinoma, kad $a$ ir $b$ yra tarpusavyje pirminiai ($DBD(a, b) = 1$)?
    \item Švyturys A mirksi kas 12 sekundžių, švyturys B – kas 15 sekundžių, o švyturys C – kas 20 sekundžių. Ką tik visi trys švyturiai sumirksėjo vienu metu. Po kiek laiko jie vėl visi trys sumirksės kartu?
    \item Įrodykite, kad skaičiai $2^n - 1$ ir $2^{n+1} - 1$ visada yra tarpusavyje pirminiai.
    \item Krumpliaratis, turintis $18$ dantukų, suka kitą krumpliaratį, turintį $45$ dantukus. Ant abiejų krumpliaračių yra po vieną raudoną atžymą. Šiuo metu atžymos liečiasi viena su kita. Kiek mažiausiai pilnų apsisukimų turi padaryti mažasis krumpliaratis, kad raudonos atžymos vėl susiliestų?
    \item Skaičius $A$ sudarytas iš lygiai $100$ vienetų ($111\dots11$), o skaičius $B$ sudarytas iš lygiai $60$ vienetų. Raskite $DBD(A, B)$.
    \item \textbf{Iššūkis:} Raskite visus natūraliuosius skaičius $n$, kuriems esant trupmena $\frac{n^2 + 4}{n + 5}$ bus \textbf{suprastinama}.
\end{enumerate}

\end{document}